%! Radicals and Rational Exponents — Unit 1, Lesson 1.2 — Warm-Up
\documentclass[10pt]{article}
\usepackage{atda-article}
\usepackage{atda-boxes}
\newcommand{\TallMath}[1]{$\displaystyle #1\rule[-1.4em]{0pt}{3.2em}$}
\setlength{\workrowsep}{6pt}   % handwriting room; moves blank and key together

\begin{document}

\pageheader{Unit 1, Lesson 1.2}{Warm-Up --- Undoing a Power}
% NAMESTRIP: no \namedateperiod here. The student writes their name once, on the
% cover this component is stapled behind. See references/conventions.md.

\vspace{0.15in}

\begin{notesbox}{Four Moves Today's Lesson Runs On}
\small
Item 4 is the one to slow down on --- it uses a rule you already know to force out a rule you
have not met yet. Show your work on all four.

\begin{enumerate}[leftmargin=1.6em, itemsep=6pt, topsep=4pt]

\item \textbf{Spiral --- Lesson 1.1.} Simplify: $\left(5y^3\right)^2$
\begin{work}
  \left(5y^3\right)^2 &= 5^2\left(y^3\right)^2 \\
                      &= 25y^6
\end{work}

\item \textbf{Perfect powers.} Evaluate each: $\sqrt{144}$, \quad $\sqrt[3]{64}$, \quad $\sqrt{81}$
\begin{work}
  \sqrt{144} &= 12 \qquad \sqrt[3]{64} = 4 \qquad \sqrt{81} = 9
\end{work}

\item \textbf{Hunting perfect squares.} Rewrite each number as (largest perfect-square factor)
$\times$ (what is left over): \quad $72$, \quad $48$, \quad $200$
\begin{work}
  72 &= 36 \cdot 2 \qquad 48 = 16 \cdot 3 \qquad 200 = 100 \cdot 2
\end{work}

\item \textbf{Discovery.} Lesson 1.1's product rule says $a^m \cdot a^n = a^{m+n}$. Apply it to
$9^{1/2} \cdot 9^{1/2}$, then finish the sentence: ``So $9^{1/2}$ must be the number that\ldots''
\begin{work}
  9^{1/2} \cdot 9^{1/2} &= 9^{\frac12+\frac12} = 9^1 = 9 \\
  \text{So } 9^{1/2} &\text{ is the number that squares to } 9, \text{ namely } 3.
\end{work}

\end{enumerate}
\end{notesbox}

\end{document}
